\section{Framework for homomorphic mapping}
\label{sec-model}


We provide \HFrame, a framework~\revise{for the subgraph homomorphic problem,}
by integrating algorithmic techniques and machine learning models.
The framework is outlined in Figure~\ref{fig:workflow}.
Given a graph $G$ and a pattern $Q$,
(1) it first leverages \kw{DualSim}
to refine the set ${\mathcal C}(u)$ of 
candidates for each vertex $u$ in $Q$ and subsequently constructs $S_{(Q, G)}$.
Recall that ${\mathcal C}(u)$ initially 
	consists of vertices in $G$ that carry the same label as $u$
	(see Section~\ref{sec-intro}).
%
(2) Then it checks whether ${\mathcal C}(u)$
is empty \revise{for some vertex $u$ in $Q$};
if so, then it returns \false, \ie 
no homomorphic mapping from $Q$ to $G$ exists;
otherwise,  
it constructs a subgraph $G'$ of $G$ 
	using only vertices appearing in $S_{(Q, G)}$; 
intuitively, it removes irrelevant vertices and edges from $G$,
which improves the performance of \HGIN (see Section~\ref{exp:sec-exp}).
(3) Finally, it  picks a {\em pivot vertex} $u_0$ from $Q$, and invokes ML model \HGIN
to check  
the existence of a homomorphic mapping  
$\varphi$ from $Q$ to $G'$
 such that $\varphi(u_0)=v$
for each 
vertex $v$ in 
${\mathcal C}(u_0)$, and 
	returns \true if \HGIN predicts 
	that such a mapping exists for some $v$ in $G'$.



  

